\documentclass[reqno]{amsart}
\usepackage{a4wide}
\usepackage{hyperref}

\AtBeginDocument{{\noindent\small
\emph{Electronic Journal of Differential Equations},
Vol. 2026 (2026), No. 11, pp. 1--10.\newline
ISSN: 1072-6691. URL: https://ejde.math.txstate.edu, DOI: 10.58997/ejde.2026.11}
\thanks{\copyright 2026. This work is licensed under a CC BY 4.0 license.}
\vspace{8mm}}

\begin{document}
\title[\hfilneg EJDE-2026/11\hfil Perturbed Schr\"odinger-Bopp-Podolsky systems]
{Existence of solutions with prescribed frequency for perturbed 
Schr\"odinger-Bopp-Podolsky systems in bounded domains}

\author[D. G. Afonso, B. Mascaro \hfil EJDE-2026/11\hfilneg]
{Danilo Gregorin Afonso, Bruno Mascaro}

\address{Danilo Gregorin Afonso\newline
Universit\`a San Raffaele Roma,
Via di Val Cannuta 247, 00166, Roma, Italy}
\email{danilo.afonso@uniroma5.it}

\address{Bruno Mascaro \newline
Faculdade de Computa\c{c}\~ao e Inform\'atica,
Universidade Presbiteriana Mackenzie,
R. da Consola\c{c}\~ao 930, S\~ao Paulo, Brasil.\newline
Instituto Mau\'a de Tecnologia, 
Pra\c{c}a Mau\'a 1, S\~ao Caetano - SP, Brasil}
\email{bruno.mascaro@mackenzie.br}

\thanks{Submitted October 30, 2025. Published February 11, 2026.}
\subjclass[2020]{35D30, 35J10, 35J20, 35J35, 35J40, 35J58, 35J91, 35Q40}
\keywords{Schr\"odinger-Bopp-Podolsky systems; Dirichlet boundary conditions;
\hfill\break\indent higher-order elliptic problems; variational methods}

\begin{abstract}
In this article, we show that the Schr\"odinger-Bopp-Podolsky system with Dirichlet
boundary conditions in a bounded domain possesses infinitely many solutions of
prescribed frequency, for any set of (continuous) boundary conditions, provided that
the Schr\"odinger equation is perturbed with a suitable nonlinearity. Our approach
is variational, and our proof is based on a symmetric variant of the Mountain Pass
theorem.
\end{abstract}

\maketitle
\numberwithin{equation}{section}
\newtheorem{theorem}{Theorem}[section]
\newtheorem{lemma}[theorem]{Lemma}
\newtheorem{proposition}[theorem]{Proposition}
\newtheorem{remark}[theorem]{Remark}
\allowdisplaybreaks

\section{Introduction} 
\label{sec:intro}

In this article, we analyze the existence of solutions to the so-called Schr\"odinger-Bopp-Podolsky system,
\begin{equation} \label{eq:pde}
  \begin{gathered}
 - \frac{1}{2} \Delta u + \phi u - g(x, u)  =  \omega u  \quad \text{in } \Omega \\
 - \Delta \phi + \Delta^2 \phi  =  4 \pi u^2  \quad \text{in } \Omega,
 \end{gathered}
\end{equation}
where $\Omega$ is a smooth bounded domain in $\mathbb R^3$, $g$ is a suitable 
nonlinearity and $\Delta^2 \phi = \Delta (\Delta \phi)$ is the bi-Laplacian operator. 
We consider Dirichlet boundary conditions, i.e.,
\begin{equation}
 \label{eq:Dirichlet_BC}
\begin{gathered}
 u  =  0  \quad \text{ on } \partial\Omega \\
 \phi  =  h_1  \quad \text{ on } \partial\Omega \\
 \Delta \phi  =  h_2  \quad \text{ on } \partial\Omega,
 \end{gathered}
\end{equation}
and assume, for simplicity, that $h_1, h_2 \in C(\partial\Omega)$.
We refer to \cite{GazzolaGrunauSweers2010} for a discussion of appropriate boundary 
operators for higher-order elliptic problems in bounded domains.

The system of equations \eqref{eq:pde} models the (stationary) interaction of a
charged particle with an electromagnetic field with the ansatz that the wave function 
is of the form
\[
 \psi(x, t) = u(x) e^{i \omega t},
\]
where $u$ plays the role of the amplitude of the wave and $\omega$ is the frequency. 
To our knowledge, the first variational analysis of this kind of system appeared in 
\cite{dAveniaSiciliano2019} (see also 
\cite{QuoirinSicilianoSilva2024, MascaroSiciliano2022}) in the case of the whole space 
$\mathbb R^3$. In fact, \eqref{eq:pde} is a refinement of the much more studied 
Schr\"odinger-Maxwell system, introduced in \cite{BenciFortunato1998} 
(see also \cite{PisaniSiciliano2007b,PisaniSiciliano2007,PisaniSiciliano2013} and the 
references therein), where the second equation is just $- \Delta \phi = 4\pi u^2$.
 Besides the physical motivation, which consists of trying to overcome the so-called 
 infinity problem of classical Maxwell theory (we refer to \cite{dAveniaSiciliano2019} 
for more on the physical aspects of the problem), the addition of the bi-Laplacian 
in the second equation gives rise to many interesting mathematical phenomena also when 
one considers boundary value problems in bounded domains, see e.g.\
\cite{AfonsoSiciliano2021}.

The main approaches developed to treat \eqref{eq:pde} are variational, and a choice is 
to be made. One can consider a normalization condition of the type
\[
 \int_\Omega u^2 \,dx = 1,
\]
in which case the parameter $\omega$ appears as a Lagrange multiplier (as done, e.g., 
in \cite{BenciFortunato1998, PisaniSiciliano2013, AfonsoSiciliano2021}), or one can 
perform a parametric analysis on the values of $\omega$ for which the free problem has 
a solution (see \cite{PisaniSiciliano2007b}).

In this work, we follow the second approach. We are able to show that, provided that we 
perturb the Schr\"odinger equation with a suitable nonlinearity, then for any prescribed 
frequency $\omega$ and any set of continuous boundary conditions $h_1$ and $h_2$, 
the system \eqref{eq:pde}-\eqref{eq:Dirichlet_BC} possesses infinitely many solutions 
$(u, \phi)$ (see Theorem \ref{thm:main_Dirichlet}).

This work is organized as follows. In Section \ref{sec:prelim} we collect the main 
notations and definitions, and recall an important result that will be useful in our 
proofs. Section \ref{sec:existence} is devoted to the variational analysis of the 
perturbed problem and the proof of our main result, Theorem \ref{thm:main_Dirichlet}. 
We address the question of non-existence for the unperturbed problem in 
Section \ref{sec:nonexistence}.

\section{Preliminaries}
\label{sec:prelim}

\subsection{Notations and definitions}

Throughout the paper, $\Omega$ is a smooth, bounded domain (connected open set). 
For $1 \leq p \leq \infty$, $\|\cdot\|_p$ denotes the $L^p$ norm (whether on $\Omega$ 
or $\partial\Omega$ will be clear from the context). As usual, we denote by $H_0^1(\Omega)$ the
completion of $C_c^\infty(\Omega)$ with respect to the Sobolev norm $W^{1, 2}(\Omega)$. 
However, we consider $H_0^1(\Omega)$ with the equivalent norm
\[
 \|u\| := \|\nabla u\|_2, \quad u \in H_0^1(\Omega).
\]
Its dual space is denoted by $H^{-1}(\Omega)$.

The eigenvalues of $- \Delta$ in $H_0^1(\Omega)$ (counted with multiplicity) are 
denoted by $\lambda_k$, with $k \in \mathbb N$. The corresponding eigenspaces are denoted 
by $H_k$.

The Sobolev space $W^{2, 2}(\Omega)$ is, as usual, denoted by $H^2(\Omega)$. 
We also consider the functional space
\[
 \mathcal{H}(\Omega) := H^2(\Omega) \cap H_0^1(\Omega)
\]
endowed with the equivalent norm
\[
 \|\varphi\|_{\mathcal{H}(\Omega)} := \|\Delta \varphi\|_2, \quad \varphi \in \mathcal{H}(\Omega).
\]

Recall that a $C^1$ functional $J$ defined in a Banach space $E$ is said to satisfy the 
Palais-Smale condition if any sequence $(u_n)_{n \in \mathbb N}$ for which 
$(J(u_n))_{n \in \mathbb N}$ is bounded and $J'(u_n) \to 0$ as $n \to \infty$ possesses 
a convergent subsequence.

In our proofs, we perform some estimates and denote by $c_j$, $j \in \mathbb N$, some 
positive constants appearing in these estimates and whose exact values are not of 
interest.

For the nonlinearity $g \in C(\overline \Omega \times \mathbb R)$ appearing in 
\eqref{eq:pde} we use the following assumptions:
\begin{itemize}
\item[(A1)] %\label{g1}
$g$ is anti-symmetric: $g(x, - \xi) = - g(x, \xi)$ for all $x \in \overline \Omega, \xi \in \mathbb R$;

\item[(A2)] % \label{g2} 
 $g$ satisfies
\[
 \lim_{\xi \to 0} \frac{g(x, \xi)}{\xi} = 0 \quad \text{uniformly in } x;
\]

\item[(A3)] %\label{g3} 
there exist constants $a_1, a_2 \geq 0$ and $p \in (4, 6)$ such that for any $x \in \overline \Omega$ and $\xi \in \mathbb R$ it holds
\[
 |g(x, \xi)| \leq a_1 + a_2 |\xi|^{p - 1};
\]

\item[(A4)] %\label{g4} 
there exist $r > 0$ and $\mu > 4$ such that for any $\xi \in \mathbb R$ with $|\xi| \geq r$ and any $x \in \overline \Omega$ it holds
 \[
 0 \leq \mu G(x, \xi) \leq \xi g(x, \xi),
 \]
where
 \[
 G(x, \xi) = \int_0^\xi g(x, t) \,dt.
 \]
\end{itemize}

\subsection{An abstract critical point theorem}

Here, we recall a version of the Mountain Pass Theorem for functionals that 
are symmetric with respect to the action of the group 
$\mathbb Z_2 = \{- \text{id}, \text{id}\}$ that will be useful in our proof.

\begin{theorem}[{\cite[Theorem 9.12]{Rabinowitz1986}}]
 \label{thm:Mountain_Pass}
Let $E$ be an infinite-dimensional Banach space, $J \in C^1(E)$ be an even 
functional such that $J(0) = 0$ and $J$ satisfies the Palais-Smale condition. 
Suppose that $E = V \oplus X$, where $V$ is finite-dimensional and $J$ satisfies
\begin{enumerate}
 \item There exist constants $\rho, \alpha > 0$ such that
 \[
 J_{B_\rho \cap X} \geq \alpha;
 \]

 \item for each finite dimensional subspace $\widetilde E$ of $E$, there exists an $R = R(\widetilde E)$ such that $J \leq 0$ in $E \setminus B_{R(\widetilde E)}$.
 \end{enumerate}
Then $I$ possesses an unbounded sequence of critical values, and therefore there exist infinitely many critical points.
\end{theorem}

\section{Existence of solutions}
\label{sec:existence}
\subsection{Variational framework}

To perform a variational analysis of the problem, it is convenient to slightly 
modify our system so as to make all boundary conditions homogeneous. 
To this aim, we consider the auxiliary problem:
\begin{equation}
 \label{eq:auxiliary_problem_Dirichlet}
\begin{gathered}
 - \Delta \chi + \Delta^2 \chi  =  0  \quad \text{in } \Omega \\
 \chi  =  h_1  \quad \text{on } \partial\Omega \\
 \Delta \chi  =  h_2  \quad \text{on } \partial\Omega.
 \end{gathered}
\end{equation}

\begin{lemma}
 \label{lem:exitence_auxiliary_problem}
Let $h_1, h_2 \in C(\Omega)$. Then there exists a weak solution 
$\chi \in H^2(\Omega)$ to \eqref{eq:auxiliary_problem_Dirichlet}. 
Moreover, $\chi \in C^4(\Omega) \cap C(\overline \Omega)$.
\end{lemma}

\begin{proof}
Notice that substituting $\theta = \Delta \chi$ we obtain
 \[
 - \Delta \chi + \Delta^2 \chi = \Delta \theta - \theta = 0 \quad \text{ in } \Omega.
 \]
Now, by well-known results of linear elliptic equations (see, e.g., \cite{Brezis2010}), 
the problem
\begin{gather*}
 - \Delta \theta + \theta  =  0  \quad \text{in } \Omega \\
 \theta  =  h_2  \quad \text{on } \partial\Omega
 \end{gather*}
admits a (unique) weak solution $\theta \in H^1(\Omega)$. 
Moreover, standard regularity estimates (see, e.g., \cite{GilbargTrudinger2001}) 
yield $\theta \in C^2(\Omega) \cap C(\overline \Omega)$.
It is also well-known that the problem
 \begin{gather*}
 - \Delta \chi  =  \theta  \quad \text{in } \Omega \\
 \chi  =  h_1 \quad \text{on } \partial\Omega
 \end{gather*}
admits a unique weak solution $\chi \in H^1(\Omega)$. Moreover, regularity theory 
yields $\chi \in C^4(\Omega) \cap C(\overline \Omega)$.
The proof is complete, since $\chi$ thus found satisfies 
\eqref{eq:auxiliary_problem_Dirichlet}.
\end{proof}

Next, we make the change of variables
\[
 \varphi = \phi - \chi.
\]
Then we write the system in the variables $(u, \varphi)$ as
\begin{equation}
 \label{eq:modified_system_Dirichlet}
\begin{gathered}
 - \frac{1}{2} \Delta u + (\varphi + \chi)u - g(x, u) 
 =  \omega u  \quad \text{in } \Omega \\
 - \Delta \varphi + \Delta^2 \varphi  =  4 \pi u^2  \quad \text{in } \Omega \\
 u  = 0 \quad \text{on } \partial\Omega \\
 \varphi  =  0  \quad \text{on } \partial\Omega \\
 \Delta \varphi  =  0  \quad \text{on } \partial\Omega.
 \end{gathered}
\end{equation}

Let us consider the functional $F_\omega: H_0^1(\Omega) \times \mathcal{H}(\Omega) \to \mathbb R$ given by
\begin{align*}
 F_\omega(u, \varphi)
 & = \frac{1}{4} \int_\Omega |\nabla u|^2 \,dx 
 + \frac{1}{2} \int_\Omega (\varphi + \chi - \omega) u^2 \,dx 
 - \int_\Omega G(x, u) \,dx  \\
 & \quad - \frac{1}{16 \pi} \int_\Omega |\nabla \varphi|^2 \,dx 
 - \frac{1}{16 \pi} \int_\Omega |\Delta \varphi|^2 \,dx.
\end{align*}
All terms of this functional, except for $\int_\Omega G(x, u) \,dx$, are linear 
or quadratic forms in the variables $u$ and $\varphi$, and thus are of class 
$C^1$ (see \cite{BadialeSerra2011}). Moreover, also $\int_\Omega G(x, u) \,dx$ 
is of class $C^1$, due to the subcritical growth assumption (A2)
(see, e.g, \cite[Appendix B]{Rabinowitz1986}). 
Therefore $F_\omega \in C^1(H_0^1(\Omega) \times \mathcal{H}(\Omega))$.
Straightforward computations yield the following expressions for the partial 
derivatives:
\begin{gather}
  \frac{\partial F_\omega}{\partial u}(u, \varphi)[v] 
 = \frac{1}{2} \int_\Omega \nabla u \nabla v \,dx 
  + \int_\Omega (\varphi + \chi - \omega) u v \,dx 
  - \int_\Omega g(x, u) v \,dx, \quad v \in H_0^1(\Omega), \label{eq:F_u} \\
  \frac{\partial F_\omega}{\partial \varphi}(u, \varphi)[\eta] = \frac{1}{2} 
  \int_\Omega \eta u^2 \,dx - \frac{1}{8\pi} \int_\Omega \nabla \varphi \nabla \eta 
  - \frac{1}{8\pi} \int_\Omega \Delta \varphi \Delta \eta \,dx, \quad \eta \in \mathcal{H}(\Omega). \label{eq:F_varphi}
\end{gather}

From the expressions of the partial derivatives, we readily obtain that

\begin{proposition}
 \label{prop:variational_formulation_Dirichlet}
 The pair $(u, \varphi) \in H_0^1(\Omega) \times \mathcal{H}(\Omega) $ is a weak solution to
\eqref{eq:modified_system_Dirichlet} if and only if $(u, \varphi)$ is a critical point 
of $F_\omega$ in $H_0^1(\Omega) \times \mathcal{H}(\Omega) $.
\end{proposition}

Observe that the functional is strongly indefinite both from above and from below. 
More precisely, for every fixed pair $(u, \varphi) \in H_0^1(\Omega) \times \mathcal{H}(\Omega) $, we
have $F_\omega(tu, \varphi) \to + \infty$ as $t \to + \infty$, because the gradient 
term ``wins" against the $L^p$-subcritical terms (due to the Sobolev embeddings). 
Similarly, $F_\omega(u, t \varphi) \to - \infty$ as $t \to + \infty$, because of the 
quadratic terms with negative sign.

Because of this fact, standard variational methods do not apply directly, as was 
already noticed in \cite{BenciFortunato1998}. Indeed, a key idea of the argument 
presented in \cite{BenciFortunato1998} (and successfully employed by many other 
authors) is to perform a suitable ``substitution" in the system, thereby reducing 
the study to an analysis of a functional of a single variable.

The procedure goes as follows. For each $u \in H_0^1(\Omega)$, we consider the 
unique weak solution $\Phi_u \in \mathcal{H}(\Omega) $ of the problem
\begin{equation}
 \label{eq:equation_for_Phi_Dirichlet}
\begin{gathered}
 - \Delta \varphi + \Delta^2 \varphi  =  4 \pi u^2  \quad \text{in } \Omega \\
 \varphi  =  0  \quad \text{on } \partial\Omega \\
 \Delta \varphi  =  0  \quad \text{on } \partial\Omega,
 \end{gathered}
\end{equation}
which can be shown to exist by a slight modification of the argument presented 
in the proof of Lemma \ref{lem:exitence_auxiliary_problem}. 
In this way, we can define a map
\[
 \Phi: u \in H_0^1(\Omega) \mapsto \Phi(u) = \Phi_u \in \mathcal{H}(\Omega) .
\]
Observe that the map $\Phi$ is implicitly defined by the equation
\begin{equation}
 \label{eq:implicit_def_Phi_Dirichlet}
 \frac{\partial F_\omega}{\partial \varphi} (u, \varphi) = 0, \quad \varphi \in \mathcal{H}(\Omega) ,
\end{equation}
which is nothing more than the weak formulation of 
\eqref{eq:equation_for_Phi_Dirichlet}.

Observe that $\Phi$ is even and $\Phi(0) = 0$. Moreover, we have the following result:

\begin{lemma}
 \label{lem:map_Phi}
The map $\Phi$ is of class $C^1$ and bounded. Moreover,
 \begin{equation}
 \label{eq:bound_Phi}
 \|\Phi(u)\| \leq C\|u\|^2
 \end{equation}
 for some $C > 0$.
\end{lemma}

\begin{proof}
To show that $\Phi$ is of class $C^1$, we use the implicit formulation 
\eqref{eq:implicit_def_Phi_Dirichlet}. Note that the derivatives of 
$\frac{\partial F_\omega}{\partial \varphi}$ are given by
\begin{equation}
\begin{gathered}
  \frac{\partial^2 F_\omega}{\partial \varphi \partial u} (u, \varphi)[\eta, v]
  = \int_\Omega u \eta v \,dx,  \\
  \frac{\partial^2 F_\omega}{\partial \varphi^2} (u, \varphi) [\eta, \zeta]
   = - \frac{1}{8\pi} \Big(\int_\Omega \nabla \eta \nabla \zeta \,dx
   + \int_\Omega \Delta \eta \Delta \zeta \,dx \Big),
\end{gathered}
\end{equation}
which are readily seen to be continuous. Therefore $\Phi$ is of class $C^1$
(see, e.g., \cite{AmbrosettiProdi1993}).

Next, we  use \eqref{eq:equation_for_Phi_Dirichlet} and the Sobolev embeddings to 
show that $\Phi$ is bounded in $H_0^1(\Omega)$. Indeed, multiplying 
\eqref{eq:equation_for_Phi_Dirichlet} by $\Phi(u)$ and integrating by parts twice we 
obtain
\[
 \int_\Omega |\nabla \Phi(u)|^2 \,dx + \int_\Omega |\Delta \Phi(u)|^2 \,dx = 4 \pi \int_\Omega \Phi(u) u^2 \,dx
\]
Therefore, by invoking well-known Sobolev embeddings, we have
\[
 \|\Phi(u)\|_{\mathcal{H}(\Omega)}^2
  \leq 4 \pi \|\Phi(u)\|_2 \|u^2\|_2  
  \leq c_1 \|\Phi(u)\|_{\mathcal{H}(\Omega)} \|u\|_4^2
  \leq c_2 \|\Phi(u)\|_{\mathcal{H}(\Omega)} \|u\|^2,
\]
 which completes the proof.
\end{proof}

With the map $\Phi$ at hand, we can define the  reduced functional
\[
 J_\omega : u \in H_0^1(\Omega) \mapsto F_\omega(u, \Phi(u)) \in \mathbb R,
\]
which can be written as
\begin{equation}
\begin{aligned}
 J_\omega(u)
 & = \frac{1}{4} \int_\Omega |\nabla u|^2 \,dx
  + \frac{1}{2} \int_\Omega (\chi - \omega) u^2 \,dx
  - \int_\Omega G(x, u) \,dx  \\
 & \quad + \frac{1}{16 \pi} \int_\Omega |\nabla \Phi(u)|^2 \,dx
 + \frac{1}{16 \pi} \int_\Omega |\Delta \Phi(u)|^2 \,dx.
\end{aligned}\label{eq:def_J}
\end{equation}
By using the chain rule together with \eqref{eq:implicit_def_Phi_Dirichlet}, we obtain
\begin{equation}
\begin{aligned}
 J_\omega'(u)[v]
 & = \frac{\partial F_\omega}{\partial u}(u, \Phi(u))[v]
  + \Big(\frac{\partial F_\omega}{\partial \varphi
 }(u, \Phi(u)) \circ \Phi'(u)\Big) [v]  \\
 & = \frac{1}{2} \int_\Omega \nabla u \nabla v \,dx + \int_\Omega (\Phi(u)
 + \chi - \omega) u v \,dx - \int_\Omega g(x, u) v \,dx.
\end{aligned} \label{eq:J_prime}
\end{equation}

The next result tells us that the problem \eqref{eq:modified_system_Dirichlet} 
can  be studied through the functional $J_\omega$.

\begin{proposition}
 \label{prop:equivalence_F_J_Dirichlet}
The pair $(u, \varphi) \in H_0^1(\Omega) \times \mathcal{H}(\Omega) $ is a critical point of 
$F_\omega$ if and only if $u$ is a critical point of $J_\omega$ and 
$\varphi = \Phi(u)$.
\end{proposition}

\begin{proof}
Suppose $(u, \varphi)$ is a critical point for $F_\omega$. 
From \eqref{eq:implicit_def_Phi_Dirichlet}, it follows that $\varphi = \Phi(u)$, 
and then it follows from \eqref{eq:F_u} that $J_\omega'(u) = 0$.

Conversely, if $\varphi = \Phi(u)$ then 
$\frac{\partial F_\omega}{\partial \varphi}(u, \Phi(u)) = 0$ by 
\eqref{eq:implicit_def_Phi_Dirichlet}, and 
$\frac{\partial F_\omega}{\partial u}(u, \Phi(u)) = 0$ since 
$J_\omega'(u) = 0$ (taking into account \eqref{eq:F_u} and \eqref{eq:J_prime}).
\end{proof}

\subsection{Analysis of the reduced functional}

We begin by showing that the reduced functional $J_\omega$ defined in \eqref{eq:def_J} 
satisfies the Palais-Smale condition.

\begin{lemma}
 \label{lem:PS_condition}
The functional $J_\omega$ defined in \eqref{eq:def_J} satisfies the Palais-Smale 
condition.
\end{lemma}

\begin{proof}
 Let $(u_n)_{n \in \mathbb N}$ be a Palais-Smale sequence, that is,
 \begin{gather}
  |J(u_n)| \leq M \quad \forall n \in \mathbb N, \text{ and  some } M > 0, 
   \label{eq:PS_a} \\
  J'(u_n) \to 0 \quad \text{in } H^{-1} \text{ as } n \to \infty. \label{eq:PS_b}
 \end{gather}
As usual in this type of argument, the first step is to show that the sequence 
$(u_n)_{n \in \mathbb N}$ is bounded. To this aim, we take $r$ as in (A4)
and  use  \eqref{eq:PS_a} to obtain
\begin{equation}
\begin{aligned}
&\frac{1}{4} \int_\Omega |\nabla u|^2 \,dx
  + \frac{1}{2} \int_\Omega (\chi - \omega) u^2 \,dx
  + \frac{1}{16 \pi} \int_\Omega |\nabla \Phi(u)|^2 \,dx
  + \frac{1}{16 \pi} \int_\Omega |\Delta \Phi(u)|^2 \,dx  \\
& \leq M + \int_{\{x \in \Omega \ : \ |u_n(x)| < r\}} |G(x, u_n)| \,dx
 + \int_{\{x \in \Omega  :  |u_n(x)| \geq r\}} |G(x, u_n)| \,dx  \\
& \leq M_1 + \frac{1}{\mu} \int_{\{x \in \Omega  : \|u_n(x)| \geq r\}} g(x, u_n) u_n \,dx  \\
& \leq M_2 + \frac{1}{\mu} \int_\Omega g(x, u_n) u_n \,dx \quad
\forall n \in \mathbb N,
\end{aligned}\label{eq:estimate_PS_a}
\end{equation}
where $M_1$ and $M_2$ are suitable positive constants.

On the other hand, since
 \begin{align*}
 |J_\omega'(u_n)[u_n]|
& = \Big|\frac{1}{2} \int_\Omega |\nabla u_n|^2 \,dx 
  + \int_\Omega (\Phi(u) + \chi - \omega) u_n^2 \,dx 
  - \int_\Omega g(x, u_n) u_n \,dx\Big|  \\
& \geq \Big|\, \Big|\int_\Omega g(x, u_n) u_n \,dx\Big| 
  - \Big|\frac{1}{2} \int_\Omega |\nabla u_n|^2 \,dx 
   + \int_\Omega (\Phi(u_n) + \chi - \omega) u_n^2 \,dx\Big|\, \Big|  \\
& \geq \int_\Omega g(x, u_n) u_n \,dx 
 - \Big|\frac{1}{2} \int_\Omega |\nabla u_n|^2 \,dx 
 + \int_\Omega (\Phi(u_n) + \chi - \omega) u_n^2 \,dx\Big|,
\end{align*}
from \eqref{eq:PS_b} we obtain that there exists some positive constant $M_3$ such 
that $|J_\omega'(u_n)[u_n]| \leq M_3 \|u_n\|$, and therefore
\begin{align*}
 \int_\Omega g(x, u_n) u_n \,dx
 & \leq M_3 \|u_n\| + \frac{1}{2} \int_\Omega |\nabla u_n|^2 \,dx 
  + \int_\Omega (\Phi(u_n) + \chi - \omega) u_n^2 \,dx  \\
 & = M_3 \|u_n\| + \frac{1}{2} \|u_n\|^2 + \int_\Omega \chi u_n^2 \,dx 
  - \omega \|u_n\|_2^2 + \frac{1}{4\pi} \|\Phi(u_n)\|^2.
\end{align*}
Substituting this into \eqref{eq:estimate_PS_a}, we obtain
 \begin{align*}
&\frac{1}{4} \int_\Omega |\nabla u|^2 \,dx
  + \frac{1}{2} \int_\Omega (\chi - \omega) u^2 \,dx 
  + \frac{1}{16 \pi} \int_\Omega |\nabla \Phi(u)|^2 \,dx 
  + \frac{1}{16 \pi} \int_\Omega |\Delta \Phi(u)|^2 \,dx  \\
& \leq M_2 + \frac{1}{\mu}\Big(M_3 \|u_n\| + \frac{1}{2} \|u_n\|^2 
 + \int_\Omega \chi u_n^2 \,dx - \omega \|u_n\|_2^2 
 + \frac{1}{4\pi} \|\Phi(u_n)\|^2 \Big)  \\
& \leq M_2 + \frac{1}{\mu}\Big(M_3 \|u_n\| + \frac{1}{2} \|u_n\|^2 
 + (\|\chi\|_\infty - \omega) \|u_n\|_2^2 + \frac{1}{4\pi} \|\Phi(u_n)\|^2 \Big) 
 \quad \forall n \in \mathbb N.
 \end{align*}
Since $\mu > 4$, we have
 \[
 \frac{1}{4\pi \mu} \|\Phi(u_n)\|^2 
 - \frac{1}{16\pi} \Big(\int_\Omega |\nabla \Phi(u_n)|^2 \,dx 
 + \int_\Omega |\Delta \Phi(u_n)|^2 \,dx \Big) < 0,
 \]
and therefore
\begin{equation}
 \label{eq:estimate_PS_b}
 \frac{\mu - 2}{4\mu} (\|u_n\|^2 - (\|\chi\|_\infty 
 - \omega)\|u_n\|_2^2) \leq M_2 + \frac{M_3}{\mu} \|u_n\| 
 \quad \forall n \in \mathbb N.
\end{equation}

If $\|\chi\|_\infty - \omega \leq 0$, we readily obtain that $(u_n)_{n \in \mathbb N}$ 
is bounded. If, instead, it holds $\|\chi\|_\infty - \omega > 0$, then from 
\eqref{eq:estimate_PS_b} we infer that
 \begin{align*}
 \|u_n\|_2^2
 & \geq \frac{c_3}{\|\chi\|_\infty - \omega} 
  \Big( \frac{\mu - 2}{4\mu}\|u_n\|^2 - \frac{M_3}{\mu}\|u_n\| - M_2 \Big)  \\
 & \geq c_4 \|u_n\|^2 - c_5\|u_n\| - c_6 \quad \forall n \in \mathbb N.
 \end{align*}
Now, should the sequence $(u_n)_{n \in \mathbb N}$ be unbounded in $H_0^1(\Omega)$, 
we would obtain $\|u_n\|_2^2 \to + \infty$ as $n \to \infty$. 
However, from assumption (A4) we deduce that there exist constants $b_1, b_2 > 0$ 
such that for any $x \in \overline \Omega$ and $\xi \in \mathbb R$ it holds
 \[
 G(x, \xi) \geq b_1 |\xi|^\mu - b_2.
 \]
But then, since $\Phi$ is bounded, we have
 \begin{align*}
 J_\omega(u_n)
& \leq \frac{1}{4}\|u_n\|^2 + \frac{1}{2}(\|\chi\|_\infty 
 - \omega) \|u_n\|_2^2 - b_1 \int_\Omega \|u_n\|^\mu \,dx 
 + b_2 |\Omega| + \frac{1}{16\pi}\|\Phi(u_n)\|^2  \\
& \leq c_7\|u_n\|^2 + c_8 \|u_n\|_2^2 - b_1 \int_\Omega |u_n|^\mu \,dx 
 + b_2|\Omega| + c_9\|u_n\|_2^2  \\
& \to - \infty \quad \text{ as } n \to \infty,
\end{align*}
since $\mu > 4$. This contradicts \eqref{eq:PS_a}, and therefore we conclude that 
$(u_n)_{n \in \mathbb N}$ is bounded also in case $\|\chi\|_\infty - \omega > 0$.

Since $(u_n)_{n \in \mathbb N}$ is bounded, then there exists $u \in H_0^1(\Omega)$ 
such that
 \[
 u_n \rightharpoonup u \quad \text{ weakly in } H_0^1(\Omega) \text{ as } n \to \infty.
 \]
We now proceed to show that the convergence is, in fact, strong in $H_0^1(\Omega)$.

To this aim, we apply \eqref{eq:J_prime} and obtain
\begin{equation}
 \label{eq:PS_J_prime}
 - \frac{1}{2} \Delta u_n = J_\omega'(u_n) - \Phi(u_n)u_n + (\chi - \omega) u_n + g(x, u_n),
\end{equation}
where this equality is understood in $H^{-1}(\Omega)$. Now, since the resolvent operator 
$(- \Delta)^{-1}: H^{-1}(\Omega) \to H_0^1(\Omega)$ is compact, to conclude our proof 
it suffices to show that the right-hand side in \eqref{eq:PS_J_prime} is bounded.

By assumption, $J_\omega'(u_n) \to 0$ as $n \to \infty$, so 
$\{J_\omega'(u_n)\}_{n \in \mathbb N}$ is bounded. Since $g$ is continuous and has 
subcritical growth, the Nemitski operator
\[
 u \in H_0^1(\Omega) \mapsto g(x, u) \in H^{-1}(\Omega)
\]
is compact (see, for example, \cite{AmbrosettiMalchiodi2007, AmbrosettiProdi1993}) and 
therefore, since $(u_n)_{n \in \mathbb N}$ is bounded, also 
$\{g(x, u_n)\}_{n \in \mathbb N}$ converges in $H^{-1}(\Omega)$. Furthermore, 
since $H_0^1(\Omega)$ is reflexive, then $H^{-1}(\Omega)$ is reflexive. 
Since $(\chi - \omega) u_n \to (\chi - \omega) u$ in the weak-$^*$ topology of 
$H^{-1}(\Omega)$, then $\{(\chi - \omega)u_n\}_{n \in \mathbb N}$ is bounded in 
$H^{-1}(\Omega)$. Finally, H\"older's inequality together with \eqref{eq:bound_Phi} 
yields
\begin{equation}
 \|\Phi(u_n) u_n\|_{3/2}^{3/2}
  \leq \|\Phi(u_n)\|_3^{3/2} \|u_n\|_3^{3/2}
  \leq c_{10} \|\Phi(u_n)\|^{3/2} \|u_n\|_3^{3/2}
  \leq c_{11} \|u_n\|^3 \|u_n\|_3^{3/2}.
\end{equation}
Hence $\{\Phi(u_n)u_n\}_{n \in \mathbb N}$ is bounded in $L^{3/2}(\Omega)$,
which is continuously embedded into $H^{-1}(\Omega)$, because $H_0^1(\Omega)$ is
continuously embedded into $L^3(\Omega)$.

Hence all terms in the right-hand side of \eqref{eq:PS_J_prime} are bounded, and 
since $(- \Delta)^{-1}$ is compact, it follows that the sequence 
$(u_n)_{n \in \mathbb N}$ converges strongly in $H_0^1(\Omega)$.
\end{proof}

\begin{proposition}
 \label{prop:Critical_points_J_Dirichlet}
For each $\omega \in \mathbb R$, the functional $J_\omega$ defined in \eqref{eq:def_J} 
has infinitely many critical points in $H_0^1(\Omega)$.
\end{proposition}

\begin{proof}
Our aim is to apply Theorem \ref{thm:Mountain_Pass} to the functional $J_\omega$. 
Since $J_\omega$ is even (because so is $\Phi$), $J_\omega(0) = 0$ and 
$J_\omega$ satisfies the Palais-Smale condition (Lemma \ref{lem:PS_condition}), 
it remains only to prove that the geometrical conditions of Theorem 
\ref{thm:Mountain_Pass} hold.

We begin by recalling from the proof of Lemma \ref{lem:PS_condition} that
\[
 J_\omega(u)
 \leq \frac{1}{4}\|u\|^2 + \frac{1}{2}(\|\chi\|_\infty - \omega) \|u\|_2^2 
 - b_1 \int_\Omega \|u\|^\mu \,dx + b_2 |\Omega| 
 + \frac{1}{16\pi}\|\Phi(u)\|^2  
 \to - \infty
\]
as $\|u\| \to + \infty$, since $\mu > 4$. Therefore, condition \textit{(ii)} 
in Theorem \ref{thm:Mountain_Pass} is satisfied.

Now, suppose that
\begin{equation}
 \label{eq:case_1_Dirichlet}
 \|\chi\|_\infty + \omega < \frac{1}{2} \lambda_1,
\end{equation}
where $\lambda_1$ is the first eigenvalue of $- \Delta$ in $H_0^1(\Omega)$ 
(see Section \ref{sec:prelim}). In this case, we can take $V = \{0\}$ and 
$X = H_0^1(\Omega)$. Indeed, $J_\omega$ has a strict local minimum at $u = 0$.
This can be proven as follows. From (A2) and (A3) we deduce that for every 
$\varepsilon > 0$ there exists $C_\varepsilon > 0$ such that for any 
$x \in \overline \Omega$ and $\xi \in \mathbb R$ it holds
\[
 |G(x, \xi)| \leq \frac{\varepsilon}{2} \xi^2 + C_\varepsilon |\xi|^p,
\]
for $p \in (4, 6)$.

Let us fix a number $c$ such that $\|\chi\|_\infty + \omega < c < \frac{1}{2} \lambda_1$. 
Since $H_0^1(\Omega) \hookrightarrow L^p(\Omega)$, by Poincar\'e inequality we obtain
\begin{align*}
 J_\omega(u)
 & \geq \frac{1}{4} \int_\Omega |\nabla u|^2 \,dx 
  - \frac{1}{2} \|\chi\|_\infty \int_\Omega u^2 \,dx 
  - \frac{1}{2} \omega \int_\Omega u^2 \,dx - \int_\Omega G(x, u) \,dx  \\
 & \geq \frac{1}{4}(\|u\|^2 - 2(\|\chi\|_\infty + \omega) \|u\|_2^2) 
  - \frac{\varepsilon}{2} \|u\|_2^2 - C_\varepsilon \|u\|_p^p  \\
 & > \frac{1}{4} \frac{\lambda_1 - c}{\lambda_1}\|u\|^2 
  - \frac{\varepsilon}{2 \lambda_1} \|u\|^2 
  - C_\varepsilon' \|u\|^p \quad \forall u \in H_0^1(\Omega),
\end{align*}
which is positive for $\|u\| = \rho$ if $\rho$ is sufficiently small. 
So condition (i) of Theorem \ref{thm:Mountain_Pass} is satisfied in the case 
when \eqref{eq:case_1_Dirichlet} holds.

Next, we assume that
\begin{equation}
 \label{eq:case_2_Dirichlet}
 \frac{1}{2} \lambda_1 \leq \|\chi\|_\infty + \omega
\end{equation}
and let $c \in \left(\frac{1}{2} \lambda_1, \|\chi\|_\infty + \omega \right)$. Set
\[
 k_\omega := \min\{k \in \mathbb N  :  \|\chi\|_\infty + \omega < \lambda_k\},
\]
and consider the following splitting 
\[
 H_0^1(\Omega) = V_k \oplus X,
\]
where
\[
 V_k = \oplus_{k = 1}^{k_\omega - 1} H_k, \quad 
 X = V_k^\perp = \overline{\oplus_{k = k_\omega}^\infty H_k}.
\]
Then, since
 \[
 \|v\|^2 \geq \lambda_{k_\omega} \|v\|^2 \quad \forall v \in X,
 \]
(by the variational characterization of the eigenvalues of $- \Delta$ in 
$H_0^1(\Omega)$), we have
\begin{align*}
 J_\omega(u)
& \geq \frac{1}{4} \int_\Omega |\nabla u|^2 \,dx 
 - \frac{1}{2} \|\chi\|_\infty \int_\Omega u^2 \,dx 
 - \frac{1}{2} \omega \int_\Omega u^2 \,dx - \int_\Omega G(x, u) \,dx  \\
& \geq \frac{1}{4}(\|u\|^2 - 2(\|\chi\|_\infty + \omega) \|u\|_2^2) 
 - \frac{\varepsilon}{2} \|u\|_2^2 - C_\varepsilon \|u\|_p^p  \\
& > \frac{1}{4} \frac{\lambda_{k_\omega} - c}{\lambda_{k_\omega}}\|u\|^2
 - \frac{\varepsilon}{2 \lambda_1} \|u\|^2 - C_\varepsilon' \|u\|^p \quad 
 \forall u \in X,
\end{align*}
so that $J_\omega$ is positive in small spheres of $X$.

We can therefore apply Theorem \ref{thm:Mountain_Pass} to conclude the proof.
\end{proof}

\begin{theorem}
 \label{thm:main_Dirichlet}
Let $\Omega \subset \mathbb R^3$ be a smooth, bounded domain. 
If $g \in C(\overline \Omega \times \mathbb R)$ satisfies {\rm (A1)--(A4)}, 
then for every triple $(h_1, h_2, \omega) \in C(\partial\Omega) \times C(\partial\Omega) \times \mathbb R$
there exist infinitely many weak solutions 
$(u, \phi) \in H_0^1(\Omega) \times H^2(\Omega)$ to the problem
 \begin{gather*}
 - \frac{1}{2} \Delta u + \phi u - g(x, u)  =  \omega u  \quad \text{in } \Omega \\
 - \Delta \phi + \Delta^2 \phi  =  4 \pi u^2  \quad \text{in } \Omega \\
 u  =  0  \quad \text{on } \partial\Omega \\
 \phi  =  h_1  \quad \text{on } \partial\Omega \\
 \Delta \phi  =  h_2  \quad \text{on } \partial\Omega.
 \end{gather*}
\end{theorem}

\begin{proof}
The statement of the theorem follows by combining Propositions 
\ref{prop:variational_formulation_Dirichlet}, 
\ref{prop:Critical_points_J_Dirichlet}, and \ref{prop:Critical_points_J_Dirichlet}.
\end{proof}

\begin{remark} \rm
 The physical meaning of Theorem \ref{thm:main_Dirichlet} is that, as happens 
in the case of Maxwell electrodynamics (\cite{PisaniSiciliano2007b}), 
if the Schr\"odinger equation is perturbed with a suitable nonlinearity, then for 
any prescribed frequency $\omega$ there exist (infinitely many) stationary solutions 
with that frequency.
\end{remark}


\section{Non-existence of solutions for the unperturbed problem}
\label{sec:nonexistence}

\begin{lemma}
 \label{lem:max_principle_Dirichlet}
Let $h_1, h_2 \in C(\partial\Omega)$ be such that $h_1 - h_2 \geq 0$ on $\partial\Omega$.
If $\phi \in H^2(\Omega)$ satisfies
 \begin{gather*}
 - \Delta \phi + \Delta^2 \phi  \geq  0  \quad \text{in } \Omega \\
 \phi  =  h_1  \quad \text{on } \partial\Omega \\
 \Delta \phi  =  h_2  \quad \text{on } \partial\Omega,
 \end{gather*}
in the weak sense, then $\phi > 0$ in $\Omega$.
\end{lemma}

\begin{proof}
Notice that 
\[
 - \Delta (\phi - \Delta \phi) \geq 0 \quad \text{ in } \Omega,
\]
so that the maximum principle yields
\[
 \phi - \Delta \phi \geq 0.
\]
Another application of the maximum principle yields the claim.
\end{proof}

\begin{proposition}
If $h_1 - h_2 \geq 0$ and $\omega < \frac{\lambda_1}{2}$, then there are no 
solutions to the problem
 \begin{gather*}
 - \frac{1}{2} \Delta u + \phi u  =  \omega u  \quad \text{in } \Omega \\
 - \Delta \phi + \Delta^2 \phi  =  4 \pi u^2  \quad \text{in } \Omega \\
 u  =  0  \quad \text{on } \partial\Omega \\
 \phi  =  h_1  \quad \text{on } \partial\Omega \\
 \Delta \phi  =  h_2  \quad \text{on } \partial\Omega .
 \end{gather*}
\end{proposition}

\begin{proof}
By Lemma \ref{lem:max_principle_Dirichlet}, $\phi > 0$ in $\Omega$. 
Multiplying the first equation by $u$ and integrating by parts, we obtain
\[
 \omega \int_\Omega u^2 \,dx > \frac{1}{2} \int_\Omega |\nabla u|^2 \,dx,
\]
which, by the variational characterization of the eigenvalues of $- \Delta$ 
in $H_0^1(\Omega)$, can only hold if $\omega > \frac{\lambda_1}{2}$.
\end{proof}

\subsection*{Acknowledgements}
This work was partially funded by the European Union - NextGenerationEU within the 
framework of PNRR Mission 4 - Component 2 - Investment 1.1 under the Italian Ministry 
of University,
and by the Research (MUR) program PRIN 2022 - Grant number 2022BCFHN2 - Advanced 
theoretical aspects in PDEs and their applications - CUP: \\ H53D23001960006.
This research was also partially supported by Gruppo Nazionale per 
l'Analisi Matematica, la Probabilit\'a e le loro Applicazioni (GNAMPA) of the 
Istituto Nazionale di Alta Matematica (INdAM).

\begin{thebibliography}{10}

\bibitem{AfonsoSiciliano2021}
Danilo Gregorin Afonso, Gaetano Siciliano;
\emph{Normalized solutions to a {S}chr\"odinger-{B}opp-{P}odolsky system under
{N}eumann boundary conditions},
 Commun. Contemp. Math. (2021).

\bibitem{AmbrosettiMalchiodi2007} Antonio Ambrosetti, Andrea Malchiodi;
\emph{Nonlinear analysis and semilinear elliptic problems}, 
Cambridge University Press, 2007.

\bibitem{AmbrosettiProdi1993} Antonio Ambrosetti, Giovanni Prodi;
\emph{A primer of nonlinear analysis},
 Cambridge University Press, 1993.

\bibitem{BadialeSerra2011} Marino Badiale, Enrico Serra;
\emph{Semilinear elliptic equations for beginners}, Springer, 2011.

\bibitem{BenciFortunato1998} Vieri Benci, Donato Fortunato;
\emph{An eigenvalue problem for the {{S}}chr\"odinger-{{M}}axwell equations}, T
 opol. Method. Nonl. An. (1998).

\bibitem{Brezis2010} Haim Brezis;
\emph{Functional Analysis, Sobolev Spaces and Partial
 Differential Equations}, Springer, New York, 2010.

\bibitem{dAveniaSiciliano2019} Pietro d'Avenia, Gaetano Siciliano;
\emph{Nonlinear Schr\"odinger equation
 in the {B}opp-{P}odolsky electrodynamics: solutions in the electrostatic
 case}, J. Differential Equations (2019).

\bibitem{GazzolaGrunauSweers2010} Filippo Gazzola, Hans-Christoph Grunau, Guido Sweers;
\emph{Polyharmonic boundary value problems}, Springer, 2010.

\bibitem{GilbargTrudinger2001} David Gilbarg, Neil S. Trudinger;
\emph{Elliptic partial differential equations of second order}, Springer, 2001.

\bibitem{MascaroSiciliano2022} Bruno Mascaro, Gaetano Siciliano;
\emph{Positive solutions for a {S}chr\"dinger-{B}opp-{P}odolsky system in $\mathbb {R}^{3}$},
Communications in Mathematics \textbf{31} (2022).

\bibitem{PisaniSiciliano2007} Lorenzo Pisani, Gaetano Siciliano;
\emph{Neumann condition in the
 {S}chr\"odinger-{M}axwell system}, Topol. Method. Nonl. An. \textbf{29}
 (2007), 251--264.

\bibitem{PisaniSiciliano2007b} Lorenzo Pisani, Gaetano Siciliano;
\emph{Note on a {S}chr\"odinger-{P}oisson system in a bounded domain},
 Appl. Math. Lett. (2007).

\bibitem{PisaniSiciliano2013} Lorenzo Pisani, Gaetano Siciliano;
\emph{{C}onstrained {S}chr\"odinger-{P}oisson {S}ystem with
 {N}on-{C}onstant {I}nteraction}, Commun. Contemp. Math. \textbf{15} (2013),
 no.~1.

\bibitem{Rabinowitz1986} Paul~H. Rabinowitz;
\emph{Minimax methods in critical point theory with
 applications to differential equations}, American Mathematical Society, 1986.

\bibitem{QuoirinSicilianoSilva2024} Humberto Ramos~Quoirin, Gaetano Siciliano, Kaye Silva;
\emph{Critical points with prescribed energy for a class of functionals depending on a
 parameter: existence, multiplicity and bifurcation results}, Nonlinearity
 \textbf{37} (2024), no.~6, 065010.

\end{thebibliography}

\end{document} 